\section{Taller resuelto: La cola vertical también produce alabeo}

\noindent\textbf{Nivel:} Intermedio.

\apartado{1. Enunciado}{Una fuerza de cola Y=-1250 N actúa en cuerpo en r=(-5.2,0,-1.1) m respecto al CG. Calcule momentos de alabeo y guiñada y explique por qué la altura de cola acopla movimientos.}

\apartado{2. Datos}{Fuerza F=(0,-1250,0) N; punto por encima del CG tiene z negativo.}

\apartado{3. Incógnitas}{Las magnitudes y comprobaciones solicitadas en el caso.}

\apartado{4. Sistema de referencia}{Cuerpo dextrógiro con z abajo; momentos respecto al CG.}

\apartado{5. Hipótesis físicas}{Fuerza puntual equivalente; se omiten otros momentos aerodinámicos de cola.}

\apartado{6. Modelo matemático}{La cola vertical también produce alabeo}

\apartado{7. Ecuaciones simbólicas}{Se parte de los balances y definiciones siguientes.}

\begin{equation}\vect M=\vect r\times\vect F\end{equation}

\begin{equation}\mathcal L=-zY,\quad\mathcal N=xY\end{equation}

\apartado{8. Desarrollo algebraico}{Las operaciones se realizan en este orden, conservando las condiciones del modelo.}

\begin{enumerate}

\item Se identifica el signo de z para una cola situada por encima del CG.

\item El producto vectorial determina ambos momentos sin memorizar una regla visual.

\item Se interpreta el momento de guiñada junto al alabeo generado por el mismo esfuerzo lateral.

\end{enumerate}

\apartado{9. Sustitución numérica}{Se usan las magnitudes del apartado de datos en SI; los ángulos se convierten a radianes antes de evaluar funciones trigonométricas. Las densidades indicadas como ISA se obtienen de la capa correspondiente.}

\noindent Momento de alabeo:\quad $1.1 \left(-1250\right)\simeq -1375\,\mathrm{N\,m}$.\par

\noindent Momento de guiñada:\quad $\left(-5.2\right) \left(-1250\right)\simeq 6500\,\mathrm{N\,m}$.\par

\apartado{10. Resultados con unidades}{Los resultados pertenecen al modelo y al intervalo de datos declarados.}

\begin{center}\begin{tabular}{lr}\toprule Magnitud&Resultado\\\midrule

Momento de alabeo & $-1375\,\mathrm{N\,m}$\\

Momento de guiñada & $6500\,\mathrm{N\,m}$\\

\bottomrule\end{tabular}\end{center}

\apartado{11. Interpretación}{Una superficie destinada a estabilidad direccional también puede modificar la dinámica de alabeo.}

\apartado{12. Comprobación independiente y alcance}{Expansión del determinante del producto vectorial y comprobación de unidades N m.}

\apartado{13. Limitaciones}{No se deriva la fuerza de un beta concreto ni se incluye la respuesta del resto del avión.}

\par\smallskip\noindent{\footnotesize\textit{Datos y solución originales. Las sustituciones registradas se reproducen con SymPy y Studium; eso no establece validación experimental ni revisión experta.}}
